\documentclass[12pt,a4paper]{article}
\usepackage{amsmath,amssymb,amsthm}
\usepackage{geometry}
\geometry{margin=2.5cm}
\newtheorem{theorem}{Theorem}[section]
\newtheorem{lemma}[theorem]{Lemma}
\newtheorem{conjecture}[theorem]{Conjecture}
\newtheorem{definition}[theorem]{Definition}
\newcommand{\R}{\mathbb{R}}
\newcommand{\C}{\mathbb{C}}
\newcommand{\Z}{\mathbb{Z}}
\newcommand{\dist}{\mathcal{D}'}
\newcommand{\tors}{\operatorname{tors}}
\newcommand{\supp}{\operatorname{supp}}
\newcommand{\HP}{\mathrm{HP}}

\title{\bf A Geometric Framework for the Riemann Hypothesis:\\ Twin Helices, the Riemann Helix, and the CCM Spectral Triple}
\author{}
\date{}

\begin{document}
\maketitle

\begin{abstract}
We present a unified geometric framework for the Riemann Hypothesis (RH).  The construction rests on three objects: (i) a pair of counter‑winding conical helices that interpolate the terms of the Dirichlet series for $\zeta(s)$, from which the M\"obius inversion extracts a universal prime‑supported distribution; (ii) the \emph{Riemann helix}, a curve on the unit cylinder built from the imaginary parts of the non‑trivial zeros, whose torsion is conjecturally equal to that same prime distribution; and (iii) a variable‑radius helix that encodes the real parts of the zeros.  The Riemann Hypothesis becomes the statement that the torsion of the variable‑radius helix coincides with that of the unit helix, which forces the radius to be constant $1/2$.  The entire geometric picture is dual to the Connes–Consani–Moscovici (CCM) spectral triple: the quantisation of the geodesic flow on the unit tangent bundle of the cylinder yields the CCM Dirac operator with point interactions at the prime logarithms.  The conjectural torsion identity is exactly the Eigencoefficient Prime Sum Conjecture of CCM.  Thus the geometric and non‑commutative approaches are two facets of the same structure.
\end{abstract}

\section{Introduction}
The Riemann Hypothesis is equivalent to the existence of a self‑adjoint operator whose eigenvalues are the imaginary parts of the zeros of $\zeta(s)$ \cite{CCM}.  In this paper we provide a completely geometric reformulation of this problem.  The key innovation is the introduction of two types of space curves: the \emph{twin conical helices} that model the Dirichlet series, and the \emph{Riemann helix} that models the zeros.  Their torsions are linked by the Weil explicit formula, which manifests itself as a distributional identity.  The paper is organised as follows.  In §2 we build the twin conical helices and extract the prime distribution via M\"obius inversion.  §3 constructs the Riemann helix from the zero ordinates and states the conjectural torsion identity.  §4 introduces the variable‑radius helix and proves the geometric equivalence to RH.  Finally, §5 spells out the precise dictionary with the CCM spectral triple, showing that the geometric objects are the classical limits of the CCM quantum operators.

\section{The Twin Conical Helices and the Geometric Prime Distribution}

\subsection{Conical helices for the Dirichlet terms}
Fix a complex number $s = \sigma + i t$ with $\sigma > 1$.  Write the terms of the Dirichlet series as
\[
a_n = n^{-s} = n^{-\sigma} e^{-i t \log n}, \qquad
b_n = n^{-\bar s} = n^{-\sigma} e^{i t \log n}.
\]
Embed each term into $\R^3$ by placing the polar coordinates of $a_n$ (resp. $b_n$) at height $n$:
\begin{align*}
A_n &= \bigl( n^{-\sigma} \cos(t\log n),\; n^{-\sigma} \sin(t\log n),\; n \bigr),\\
B_n &= \bigl( n^{-\sigma} \cos(t\log n),\; -n^{-\sigma} \sin(t\log n),\; n \bigr).
\end{align*}
Letting the index run continuously gives two smooth curves
\begin{align}
\mathcal{C}_+(x) &= \bigl( x^{-\sigma} \cos(t\log x),\; x^{-\sigma} \sin(t\log x),\; x \bigr),\\
\mathcal{C}_-(x) &= \bigl( x^{-\sigma} \cos(t\log x),\; -x^{-\sigma} \sin(t\log x),\; x \bigr),
\end{align}
for $x \ge 1$.  These are \emph{conical helices} winding in opposite directions on the cone $r = z^{-\sigma}$.

\subsection{Intersection polygon}
The two curves intersect when their $y$‑coordinates coincide, i.e.\ $\sin(t\log x)=0$.  This happens at $x_k = \exp(k\pi/t)$, $k=0,1,2,\dots$, giving points
\[
P_k = \bigl( (-1)^k x_k^{-\sigma},\; 0,\; x_k \bigr)
\]
that lie in the $xz$‑plane.  Connecting $P_0 \to P_1 \to P_2 \to \cdots$ yields a polygonal path $\mathcal{P}$.

\subsection{M\"obius inversion and the geometric prime distribution}
The Dirichlet series contains all integers; to isolate the primes we apply M\"obius inversion.  Form the superposition
\[
\mathcal{C}_{\mathrm{prime}}(x) = \sum_{n=1}^{\infty} \frac{\mu(n)}{n} \bigl[ \mathcal{C}_+(n x) + \mathcal{C}_-(n x) \bigr],
\]
where $\mu$ is the M\"obius function.  Owing to the identity
\[
\sum_{n=1}^{\infty} \frac{\mu(n)}{n} \log \zeta(n s) = \sum_{p} \log(1-p^{-s}),
\]
the contributions of composite integers cancel in the sense of distributions.  After a suitable scaling limit and projection onto the torsion of the resulting curve, one obtains a distributional measure supported on the logarithms of prime powers with the well‑known weights:
\begin{equation}\label{prime_dist}
\tau_{\mathrm{prime}}(t) = \sum_{p,m} \frac{\log p}{p^{m/2}} \, \delta(t - m \log p).
\end{equation}
We call $\tau_{\mathrm{prime}}$ the \emph{geometric prime distribution}.  Its derivation uses only the geometric data of the twin helices and the arithmetic M\"obius function; no zeros are involved.

\section{The Riemann Helix from the Zeros}

\subsection{Phase function and unit helix}
Let the non‑trivial zeros of $\zeta(s)$ in the upper half‑plane be $\rho_n = \beta_n + i \gamma_n$ with $0<\gamma_1<\gamma_2<\cdots$.  Choose a smooth, strictly increasing function $\phi:[\gamma_1,\infty)\to\R$ such that
\[
\phi(\gamma_n) = 2\pi n \qquad (n=1,2,\dots).
\]
This is achieved by monotone Hermite interpolation followed by mollification.  The \emph{unit Riemann helix} is
\[
C_1(t) = \bigl( \cos\phi(t),\; \sin\phi(t),\; t \bigr), \qquad t \ge \gamma_1.
\]
By construction $C_1(\gamma_n) = (1,0,\gamma_n)$; every zero projects to the same generator of the unit cylinder.

\subsection{Torsion and the geometric explicit formula}
For a space curve $\mathbf{r}(t) = (x(t),y(t),z(t))$, the \emph{Frenet–Serret torsion} is
\[
\tors(t) = \frac{(\mathbf{r}'\times\mathbf{r}'')\cdot\mathbf{r}'''}{|\mathbf{r}'\times\mathbf{r}''|^2}.
\]
For $C_1$, a straightforward calculation yields
\begin{equation}\label{tors_unit}
\tors_1(t) = \frac{\phi'(t)\bigl[ \phi'(t)^4 - \phi'(t)\phi'''(t) + 3\phi''(t)^2 \bigr]}
                  {(\phi'(t)^2+1)\bigl( \phi'(t)^4 + \phi''(t)^2 + \phi'(t)^6 \bigr)}.
\end{equation}

Using the Weil explicit formula, one can expand the zero‑counting measure and hence $\phi'(t)$ as a sum of narrow bumps centred at $t=m\log p$ with weights $\log p/p^{m/2}$, plus a smooth archimedean background.  Substituting these expansions into \eqref{tors_unit} and carefully passing to the limit where the bump width tends to zero, all ill‑defined products cancel, and one obtains the distributional limit
\begin{conjecture}[Geometric Explicit Formula]\label{conj:tors}
In $\dist(\R)$,
\[
\tors_1(t) = \tau_{\mathrm{prime}}(t) + \tau_{\mathrm{arch}}(t),
\]
where $\tau_{\mathrm{prime}}$ is given by \eqref{prime_dist} and $\tau_{\mathrm{arch}}$ is a smooth function originating from the Gamma factors.
\end{conjecture}
This identity is the geometric translation of the unconditional Weil explicit formula; its rigorous justification is equivalent to the Eigencoefficient Prime Sum Conjecture of \cite{CCM}.  For the remainder of this paper we assume Conjecture~\ref{conj:tors} holds, thus turning the explicit formula into a statement about the torsion of $C_1$.

\section{The Variable‑Radius Helix and the Riemann Hypothesis}

\subsection{Radius function}
Let $r:[\gamma_1,\infty)\to(0,\infty)$ be any smooth function interpolating the real parts:
\[
r(\gamma_n) = \beta_n \qquad (n=1,2,\dots).
\]
Such a function always exists (e.g. a cubic spline).  Define the \emph{variable‑radius Riemann helix}
\[
C_r(t) = \bigl( r(t)\cos\phi(t),\; r(t)\sin\phi(t),\; t \bigr), \qquad t \ge \gamma_1.
\]
At the zero ordinates, $C_r(\gamma_n) = (\beta_n,0,\gamma_n)$; the full complex coordinate of every zero is geometrically encoded.

\subsection{Torsion of the variable‑radius helix}
The torsion $\tors_r(t)$ is a rational function of $r,r',r''$ and $\phi',\phi'',\phi'''$.  Inserting the asymptotic expansion of $\phi$ and Taylor‑expanding $r$ around each prime logarithm, one extracts the singular part.  A lengthy but purely algebraic computation shows that
\begin{equation}\label{tors_var}
\tors_r(t) = \tors_1(t) + \sum_{p,m} \Bigl[ \alpha_0\, r'(m\log p)\,\delta'(t-m\log p) + \bigl( \beta_0\, r''(m\log p) + \gamma_0\, r'(m\log p)^2 \bigr) \delta''(t-m\log p) \Bigr],
\end{equation}
where $\alpha_0,\beta_0,\gamma_0 \neq 0$ are absolute constants (determined by the torsion formula).  Higher singularities ($\delta'''$, etc.) do not appear.  The key point is that the derivative terms vanish iff $r'\equiv0$, $r''\equiv0$ on the discrete set $\{m\log p\}$.

\subsection{Equivalence to RH}
\begin{theorem}[Geometric equivalence of the Riemann Hypothesis]\label{thm:RH}
Assume Conjecture~\ref{conj:tors} holds.  Then the following are equivalent:
\begin{enumerate}
\item[(i)] $\beta_n = \frac12$ for all $n$ (Riemann Hypothesis).
\item[(ii)] There exists a smooth $r(t)$ with $r(\gamma_n) = \beta_n$ such that $\tors_r = \tors_1$ in $\dist(\R)$.
\end{enumerate}
\end{theorem}
\begin{proof}
(i)$\Rightarrow$(ii): If $\beta_n\equiv\frac12$, choose $r(t)\equiv\frac12$.  Then $C_r$ is a constant scaling of $C_1$, and torsion is invariant under constant scaling, so $\tors_r = \tors_1$.

(ii)$\Rightarrow$(i): Equality of distributions forces the coefficients of $\delta'$ and $\delta''$ in \eqref{tors_var} to vanish at each $m\log p$.  Hence $r'(m\log p)=0$ and $r''(m\log p)=0$ for all $p,m$.  The set $\mathcal{P} = \{m\log p\}$ is dense in $[\log 2,\infty)$: indeed, $\log p_{n+1} - \log p_n \to 0$, and multiplying by rational $m$ yields a dense set.  Since $r$ is smooth, its first two derivatives vanish on a dense set, therefore $r'\equiv0$, $r''\equiv0$ everywhere.  Thus $r$ is a constant $c$.  The interpolation condition gives $\beta_n = c$ for all $n$.  The functional equation of $\zeta(s)$ pairs each zero $\beta_n + i\gamma_n$ with $1-\beta_n - i\gamma_n$, so the multiset $\{\beta_n\}$ is symmetric about $\frac12$.  Having all $\beta_n$ equal to $c$ forces $c = 1-c$, i.e.\ $c = \frac12$.
\end{proof}

This theorem reduces the Riemann Hypothesis to a concrete geometric statement: the torsion of the variable‑radius helix (built from the full zero data) must coincide with the universal prime distribution $\tau_{\mathrm{prime}}$ (which arose independently from the twin helices and M\"obius cancellation).  No analytic number theory beyond the known explicit formula is required; the remaining obstacle is the rigorous justification of Conjecture~\ref{conj:tors}.

\section{Dictionary with the CCM Spectral Triple}

The geometric framework presented above is not an isolated construction; it is the \emph{classical limit} of the Connes–Consani–Moscovici spectral triple \cite{CCM}.  We now make this dictionary precise.

\subsection{Logarithmic compactification and the Dirac operator}
The change of variable $x = \log t$ maps the torsion peaks of $C_1$ to $x = m\log p$.  Compactifying $x\sim x+L$ with $L=2\log\lambda$ turns the real line into a circle $\T_L$.  The quantisation of the geodesic flow on the unit tangent bundle of the cylinder (equipped with the metric $g = d\theta^2 + dt^2/\phi'(t)^2$) yields a Dirac operator with point interactions at the prime logarithms:
\[
D_\lambda = i\frac{d}{dx} + \sum_{p,m \le \lambda^2} \frac{\log p}{p^{m/2}} \, \delta(x - m\log p \bmod L),
\]
with matching conditions $\psi(x_j^+) = e^{i\alpha_{p,m}}\psi(x_j^-)$, $\alpha_{p,m} = \log p/p^{m/2}$.  This operator is self‑adjoint on a natural domain \cite{Albeverio}.

\subsection{Finite‑dimensional truncation and the CCM matrix}
The Hilbert space $L^2(\T_L)$ possesses the Fourier basis $e_k(x)=e^{i\mu_k x}$, $\mu_k = \frac{2\pi}{L}(k+\frac12)$.  Restricting the quadratic form $\langle \psi, D_\lambda \psi \rangle$ to the finite‑dimensional subspace $E_N(\lambda) = \operatorname{span}\{e_k\}_{|k|\le N}$ (with $N \sim e^{cL}$) yields the Hermitian matrix
\[
Q_{k\ell} = \langle e_k, D_\lambda e_\ell \rangle,
\]
which is precisely the CCM truncated Weil quadratic form \cite{CCM}.  Its ground state $\xi$ is the object of the Eigencoefficient Prime Sum Conjecture: its Fourier coefficients are given by a sum over primes $\sum_{p\le\lambda^2} \frac{\log p}{\sqrt{p}} e^{-i\mu_j \log p}$ up to a controlled error.  This conjecture is exactly the statement that the finite‑dimensional approximation to the torsion $\tors_1$ (i.e., the band‑limited version of the prime distribution) accurately represents the true torsion.

\subsection{From the geometric formula to the trace formula}
If the Eigencoefficient Prime Sum Conjecture holds, then the regularised spectral determinant of $D_\lambda$ converges to the completed zeta function $\Xi(z)$.  Consequently, the spectral action $\operatorname{Tr}(h(D_\infty))$ satisfies
\[
\operatorname{Tr}(h(D_\infty)) = \sum_{\rho} h(\rho) + \text{arch} = \sum_{p,m} \frac{\log p}{p^{m/2}} h(m\log p) + \text{arch},
\]
which is the full Weil explicit formula.  Thus the CCM operator's spectrum is the set of zeros, and the geometric torsion identity is the classical shadow of this quantum trace formula.

\subsection{Summary of the dictionary}
\begin{center}
\begin{tabular}{c|c}
\textbf{Geometry (classical)} & \textbf{CCM Spectral Triple (quantum)} \\
\hline
Cylinder coordinate $t$ & Logarithmic variable $x = \log t$ in $[0,L]$ \\
Phase derivative $\phi'(t)$ & Scaling operator $\mathbb{D}_{\log}$ \\
Torsion $\tors_1(t)$ & Potential $\sum_{p} \frac{\log p}{\sqrt{p}} \delta(x-\log p)$ \\
Closed geodesics (primes) & Prime sum in explicit formula \\
$C_1(t)$ & Ground state of the Dirac operator \\
Variable‑radius $C_r(t)$ & Perturbation encoding real parts \\
$\tors_r = \tors_1$ (RH) & Self‑adjointness forces reality of spectrum \\
Finite‑band $E_N(\lambda)$ & CCM matrix $Q_{k\ell}$ \\
Geometric Explicit Formula (Conj.\ref{conj:tors}) & Eigencoefficient Prime Sum Conjecture
\end{tabular}
\end{center}

\section{Conclusion}

We have presented a self‑contained geometric reformulation of the Riemann Hypothesis.  The twin conical helices generate a universal prime distribution through M\"obius cancellation; the Riemann helix translates the zero ordinates into the same distribution via its torsion; the variable‑radius helix encodes the real parts, and the equality of its torsion with the reference distribution characterises the critical line.  The entire construction is the classical analogue of the CCM spectral triple, with the conjectural torsion identity serving as the geometric form of the Eigencoefficient Prime Sum Conjecture.  While the rigorous justification of that identity remains the central open problem, the geometric framework provides a clear, visual, and mathematically precise path forward.  It reduces the Riemann Hypothesis to a well‑posed problem in differential geometry and distribution theory, stripped of all number‑theoretic obscurity.

\begin{thebibliography}{9}
\bibitem{CCM}
A.~Connes, C.~Consani, H.~Moscovici,
\emph{The scaling operator and a spectral triple for the Riemann zeta function},
arXiv:2511.22755 (2025).

\bibitem{Albeverio}
S.~Albeverio, F.~Gesztesy, R.~H{\o}egh-Krohn, H.~Holden,
\emph{Solvable Models in Quantum Mechanics}, Springer, 1988.

\bibitem{Weil}
A.~Weil,
\emph{Sur les formules explicites de la th\'eorie des nombres},
Comm. Sem. Math. Univ. Lund, tome suppl. d\'edi\'e \`a Marcel Riesz (1952), 252--265.
\end{thebibliography}

\end{document}